\section{Introduction}
\label{sec:intro}

Filaments are field-aligned structures of excess density and temperature that form near the last closed flux surface (LCFS) of a tokamak and move radially outward through the scrape-off layer (SOL), carrying a large part of the cross-field particle and heat flux in the SOL \citep{krasheninnikov2001scrape, zweben2007edge, dippolito2011convective}. Their radial velocity sets how far they travel before they drain to the divertor targets and thereby the load on the first wall, and analytic models give this velocity as a function of filament size and of the parameters that control how the polarization charge of the filament is dissipated. When the charge drains through the sheath at the target, the velocity falls with size as $v \propto a^{-2}$ \citep{krasheninnikov2001scrape}, and when ion polarization current closes the circuit inside the filament, the velocity rises with size as $v \propto a^{1/2}$ \citep{garcia2006fluctuations}. The two-region model of \citet{myra2006collisionality} joins these limits and predicts a velocity maximum at a critical size $a^*$ set by the sound speed, the connection length, and the collisionality.

Testing these predictions requires velocities and sizes for many filaments under known plasma conditions, and on the Mega Amp Spherical Tokamak (MAST) the main instrument for this purpose has been a Photron SA1.1 fast camera that views the outboard midplane in visible light. Earlier MAST studies used wide-angle views, identified filaments by hand or with detection software built for that view, and worked from a small number of discharges \citep{dudson2008experiments, benayed2009interelm, kirk2016filaments, farley2019identification}, each requiring a camera calibration model fitted to the view in use.

The FAIR-MAST archive \citep{jackson2024fairmast} makes the full MAST data set public and includes a recording mode of the midplane camera that earlier work did not use at scale, a $48 \times 256$ pixel strip through the outboard edge recorded at 100{,}000 frames per second. The strip resolves radial motion at 10~$\mu$s intervals but spans only 48 rows, so it constrains radial motion and radial size and little else, and because the archive includes no spatial calibration for these recordings, the pixel scale has to be recovered from the data.

Here those recordings are turned into filament statistics with an automatic pipeline built from standard computer-vision operations, namely a running-minimum background model, robust thresholding, connected-component detection, Hungarian frame-to-frame assignment, least-squares track fitting, and Lucas--Kanade optical flow as an independent velocity estimate. The pixel scale is calibrated against the equilibrium reconstruction, since the image column of the plasma limb moves with the EFIT outboard LCFS radius from shot to shot. A robust regression over a campaign then gives the scale in millimetres per pixel with a bootstrap confidence interval for each camera configuration, using only the images and EFIT, and the pipeline is run on every usable strip recording in the archive.

Applied to the archive, the pipeline yields radial velocity and radial width for \nTracksSel{} filament tracks in \nShotsStats{} MAST discharges, with a median radial velocity of \vMedAll{}~m/s, a median radial full width at half maximum of \dMedAll{}~cm, and \outwardMed\% of tracks moving outward. The pooled velocity--size slope is $\vdSlopePooled$ in log--log units and $\vdSlopePooledBlur$ after the widths are corrected for exposure blur, so the velocity is nearly independent of size over the measured range, and in two-region units the filaments sit at $\hat a \approx \aHatMin$--$\aHatMax$, where the model places its velocity maximum, although their normalized velocity $\hat v \approx \vHatMed$ is far below the model scale. Robust regression gives the dependences of the per-shot median velocity and width on plasma current, Greenwald fraction, and beam power with confidence intervals, and code and derived data accompany the paper.
